%! Rates of Change in Polar Functions — Unit 7, Lesson 7.8 — Group Activity
\documentclass[10pt]{article}
\usepackage{precalculus-article}
\usepackage{precalculus-boxes}
\newcommand{\TallMath}[1]{$\displaystyle #1\rule[-1.4em]{0pt}{3.2em}$}
\newcommand{\CourseName}{Precalculus}
\newcommand{\SchoolYear}{2026--2027}
\newcommand{\MeetingLength}{55 minutes}
\newcommand{\UnitNumberName}{Unit 7: Analytic Trigonometry and Polar Coordinates \quad}
\newcommand{\LessonNumberName}{Lesson 7.8: Rates of Change in Polar Functions}

\begin{document}

\pageheader{Unit 7, Lesson 7.8}{Group Activity}
\namepartnerperiod

\vspace{0.10in}

%----------------------------------------------------------------------
% Context
%----------------------------------------------------------------------
\begin{scenariobox}[Context: Radar Signal Strength]{sky}
A radar system sweeps an antenna in a circle. The \textbf{angle} $\theta$ is the
current direction of the sweep (in radians), and $r = f(\theta)$ represents the
\textbf{signal strength} (in arbitrary units) detected at that angle. As the antenna
sweeps, $r$ can increase (stronger signal), decrease (weaker signal), or even
go negative (a phase-inverted return).

Analyzing how $r$ changes with $\theta$ tells engineers how quickly the signal
is strengthening or weakening as the antenna rotates. Choose \textbf{one tier}
based on your readiness. Be prepared to explain your reasoning to the class.
\end{scenariobox}

\vspace{0.08in}

%----------------------------------------------------------------------
% Tier R
%----------------------------------------------------------------------
\begin{tcolorbox}[colback=white, colframe=black!40,
    title=\textbf{Tier R --- Reinforce},
    fonttitle=\bfseries, arc=2mm, left=3mm, right=3mm, top=2mm, bottom=2mm]

\textbf{Directions:} The table below gives values of $r = 2\cos\theta$ at nine angles.
Complete Parts A and B.

\vspace{6pt}
\renewcommand{\arraystretch}{1.5}
\begin{tabular}{|c|c|c|c|c|c|c|c|c|c|}
\hline
$\theta$ & $0$ & $\pi/4$ & $\pi/2$ & $3\pi/4$ & $\pi$ & $5\pi/4$ & $3\pi/2$ & $7\pi/4$ & $2\pi$ \\
\hline
$r$ & $2$ & $\approx 1.41$ & $0$ & $\approx -1.41$ & $-2$ & $\approx -1.41$ & $0$ & $\approx 1.41$ & $2$ \\
\hline
\end{tabular}

\vspace{8pt}
\textbf{Part A --- Classify each interval.}
Write P\&I (positive and increasing), P\&D (positive and decreasing),
N\&I (negative and increasing), or N\&D (negative and decreasing)
in the table below.

\vspace{4pt}
\renewcommand{\arraystretch}{1.6}
\begin{tabular}{|c|p{4.5cm}|}
\hline
\textbf{Interval} & \textbf{Classification} \\
\hline
$[0, \pi/4]$ & \underline{\hspace{4cm}} \\
\hline
$[\pi/4, \pi/2]$ & \underline{\hspace{4cm}} \\
\hline
$[\pi/2, 3\pi/4]$ & \underline{\hspace{4cm}} \\
\hline
$[3\pi/4, \pi]$ & \underline{\hspace{4cm}} \\
\hline
$[\pi, 5\pi/4]$ & \underline{\hspace{4cm}} \\
\hline
$[5\pi/4, 3\pi/2]$ & \underline{\hspace{4cm}} \\
\hline
$[3\pi/2, 7\pi/4]$ & \underline{\hspace{4cm}} \\
\hline
$[7\pi/4, 2\pi]$ & \underline{\hspace{4cm}} \\
\hline
\end{tabular}

\vspace{8pt}
\textbf{Part B --- Average rates of change.}

\vspace{6pt}
\textbf{1.}\quad Average rate of change of $r = 2\cos\theta$ from $\theta = 0$ to $\theta = \pi/2$:

\vspace{4pt}
ARC $= \dfrac{r(\pi/2) - r(0)}{\pi/2 - 0} = \dfrac{\rule{1.5cm}{0.4pt} - \rule{1.5cm}{0.4pt}}{\pi/2} = $ \underline{\hspace{3cm}}

\vspace{8pt}
\textbf{2.}\quad Average rate of change from $\theta = \pi/2$ to $\theta = \pi$:

\vspace{4pt}
ARC $= \dfrac{r(\pi) - r(\pi/2)}{\pi/2} = \dfrac{\rule{1.5cm}{0.4pt} - \rule{1.5cm}{0.4pt}}{\pi/2} = $ \underline{\hspace{3cm}}

\vspace{8pt}
\textbf{3.}\quad For the interval $[0, \pi/2]$: in the \textit{radar context}, what does the
sign of your ARC value tell you about the signal strength?

\writelines{2}

\end{tcolorbox}

\vspace{0.08in}

%----------------------------------------------------------------------
% Tier A
%----------------------------------------------------------------------
\begin{tcolorbox}[colback=white, colframe=black!40,
    title=\textbf{Tier A --- Apply},
    fonttitle=\bfseries, arc=2mm, left=3mm, right=3mm, top=2mm, bottom=2mm]

\textbf{Directions:} Use the two tables below to answer all parts.

\textbf{Table 1:} $r = 2\cos\theta$ \hfill (same as Tier R)

\vspace{4pt}
\renewcommand{\arraystretch}{1.4}
\begin{tabular}{|c|c|c|c|c|c|c|c|c|c|}
\hline
$\theta$ & $0$ & $\pi/4$ & $\pi/2$ & $3\pi/4$ & $\pi$ & $5\pi/4$ & $3\pi/2$ & $7\pi/4$ & $2\pi$ \\
\hline
$r$ & $2$ & $\approx 1.41$ & $0$ & $\approx -1.41$ & $-2$ & $\approx -1.41$ & $0$ & $\approx 1.41$ & $2$ \\
\hline
\end{tabular}

\vspace{8pt}
\textbf{Table 2:} $r = 1 - 2\sin\theta$

\vspace{4pt}
\begin{tabular}{|c|c|c|c|c|c|c|c|c|c|}
\hline
$\theta$ & $0$ & $\pi/4$ & $\pi/2$ & $3\pi/4$ & $\pi$ & $5\pi/4$ & $3\pi/2$ & $7\pi/4$ & $2\pi$ \\
\hline
$r$ & $1$ & $\approx -0.41$ & $-1$ & $\approx -0.41$ & $1$ & $\approx 2.41$ & $3$ & $\approx 2.41$ & $1$ \\
\hline
\end{tabular}

\vspace{8pt}
\textbf{1.}\quad Identify all intervals in $[0, 2\pi]$ where $r < 0$ for each function.

\vspace{4pt}
$r = 2\cos\theta$: \underline{\hspace{5cm}}

\vspace{4pt}
$r = 1 - 2\sin\theta$: \underline{\hspace{5cm}}

\vspace{8pt}
\textbf{2.}\quad Compute the average rate of change of each function over the given interval.

\vspace{4pt}
(a)\quad $r = 2\cos\theta$ from $\theta = 3\pi/4$ to $\theta = \pi$:
\quad ARC $= $ \underline{\hspace{4cm}}

\vspace{4pt}
(b)\quad $r = 2\cos\theta$ from $\theta = \pi$ to $\theta = 5\pi/4$:
\quad ARC $= $ \underline{\hspace{4cm}}

\vspace{4pt}
(c)\quad $r = 1 - 2\sin\theta$ from $\theta = \pi/4$ to $\theta = \pi/2$:
\quad ARC $= $ \underline{\hspace{4cm}}

\vspace{4pt}
(d)\quad $r = 1 - 2\sin\theta$ from $\theta = \pi/2$ to $\theta = 3\pi/4$:
\quad ARC $= $ \underline{\hspace{4cm}}

\vspace{8pt}
\textbf{3.}\quad Interpret one of your answers from Part~2 in the context of the radar scenario.
State which function and interval you are interpreting, and write a complete sentence.

\writelines{3}

\end{tcolorbox}

\vspace{0.08in}

%----------------------------------------------------------------------
% Tier E
%----------------------------------------------------------------------
\begin{tcolorbox}[colback=white, colframe=black!40,
    title=\textbf{Tier E --- Extend},
    fonttitle=\bfseries, arc=2mm, left=3mm, right=3mm, top=2mm, bottom=2mm]

\textbf{Directions:} The four-petal rose $r = \sin(2\theta)$ produces four petals.
Use the provided table (in rectangular form) to complete all parts.

\vspace{6pt}
\textbf{Table:} $r = \sin(2\theta)$

\vspace{4pt}
\renewcommand{\arraystretch}{1.4}
\begin{tabular}{|c|c|c|c|c|c|c|c|c|c|}
\hline
$\theta$ & $0$ & $\pi/4$ & $\pi/2$ & $3\pi/4$ & $\pi$ & $5\pi/4$ & $3\pi/2$ & $7\pi/4$ & $2\pi$ \\
\hline
$r$ & $0$ & $1$ & $0$ & $-1$ & $0$ & $1$ & $0$ & $-1$ & $0$ \\
\hline
\end{tabular}

\vspace{8pt}
\textbf{Part A.}\quad Classify all eight intervals in $[0, 2\pi]$ by sign and direction
(P\&I, P\&D, N\&I, N\&D, or zero). A table format is recommended.

\writelines{6}

\vspace{6pt}
\textbf{Part B.}\quad $r = \sin(2\theta)$ produces four petals. Using your classifications,
identify the interval during which each petal is traced. (Recall: a petal is traced
when $r$ starts at 0, moves away from the origin, then returns to 0.)

\vspace{4pt}
Petal 1 (Quadrant I, $\theta \in (0, \pi/2)$): traced on \underline{\hspace{3.5cm}}

\vspace{4pt}
Petal 2 (Quadrant IV, $\theta \in (\pi/2, \pi)$): traced on \underline{\hspace{3.5cm}}

\vspace{4pt}
Petal 3 (Quadrant III, $\theta \in (\pi, 3\pi/2)$): traced on \underline{\hspace{3.5cm}}

\vspace{4pt}
Petal 4 (Quadrant II, $\theta \in (3\pi/2, 2\pi)$): traced on \underline{\hspace{3.5cm}}

\vspace{8pt}
\textbf{Part C.}\quad Write a paragraph (3--5 sentences) explaining:
\begin{itemize}[leftmargin=1.4em, itemsep=2pt]
  \item How the sign of $r$ determines which ``direction'' (which quadrant)
    each petal is traced.
  \item Why Petal 2 and Petal 4 appear in the opposite quadrant from where
    the angle $\theta$ is pointing.
  \item How the P\&I / P\&D / N\&I / N\&D classification tells you whether
    the curve is moving outward or inward during petal formation.
\end{itemize}

\writelines{6}

\end{tcolorbox}

\end{document}
